\uuid{B79E}
\titre{Matrice d'une application linéaire/Inversibilité/Espace des fonctions}
\auteur{hanani ; mzari}
\organisation{exo7}
\licence{CC BY-NC-SA 4.0}
\datecreate{2020-01-20}
\niveau{L1}
\module{Algèbre}
\chapitre{Matrice}
\sousChapitre{Matrice et application linéaire}
\difficulte{3}
\theme{Matrice d'une application linéaire, Inversibilité, Espace des fonctions}
\exo7qcmid{lille-2.matrice-application.024}

\contenu{
\begin{question}
\qclassification{108.03}
\qtags{motcle=Matrice d'une application linéaire/Inversibilité/Espace des fonctions}
On considère ${\cal F}$ l'espace vectoriel des fonctions réelles engendré par les fonctions $f_1, \; f_2$ et  $f_3$ définies par  :
$f_1(x)=1, \; f_2(x)=e^x$ et  $f_3(x)=xe^x$. Soit $\phi$ l'application linéaire définie :
$$\begin{array}{rccc}\phi:& {\cal F}&\to& {\cal F}\\
& f&\to &f+f'-f'',  \end{array}$$
où $f'$ (resp. $f''$) est la dérivée première (resp. seconde) de $f$. On notera $M$ la matrice de $\phi$ dans la base ${\cal B}=\{f_1, f_2,f_3\}$. Quelles sont les assertions vraies ?
\begin{answers}
\good{$M=\left(\begin{array}{rcc}
1&0&0\\0&1&-1\\
0&0&1\\\end{array}\right)$.}
\bad{Le rang de la matrice $M$ est $2$}.
\good{$\phi $ est bijective.}
\good{$M$ est inversible et $M^{-1} =
\left(\begin{array}{rcc}1&0&0\\0&1&1\\ 0&0&1\end{array}\right)$.}
\end{answers}
\vskip2mm
\begin{explanations} La matrice de $\phi$ d'une base $ {\cal {B}}=(u_j)$ dans une base $ {\cal {B}}'=(v_i)$ est la matrice  dont la $j$ième colonne est constituée des coordonnées de $\phi(u_j)$ dans la base $ {\cal {B}}'$. Donc
$$M= \left(\begin{array}{rcc}
1&0&0\\0&1&-1\\ 0&0&1\end{array}\right).$$
Le rang d'une matrice est le nombre maximum de vecteurs colonnes ou lignes linéairement indépendants. Donc le rang de $M$ est $3$. Par conséquent, $\phi$ est bijective et $M$ est inversible. On vérifie que
$M^{-1} = \left(\begin{array}{rcc}1&0&0\\
0&1&1\\ 0&0&1\end{array}\right)$.
\end{explanations}
\end{question}
}
